% Part: sets-functions-relations
% Chapter: sets
% Section: important-sets

\documentclass[../../../include/open-logic-section]{subfiles}

\begin{document}

\olfileid{sfr}{set}{imp}
\olsection{કેટલાક મહત્ત્વના ગણો}

\begin{ex}
આપણે મુખ્યત્વે એવા ગણોનો અભ્યાસ કરીશું જેમના !!{element}s ગાણિતિક
વસ્તુઓ હોય. આવા ચાર ગણો એટલા મહત્ત્વના છે કે તેમને ખાસ નામ આપેલાં છે:
\begin{multline*}
    \Nat = \{0, 1, 2, 3, \ldots\} \\
    \shoveright{\text{પ્રાકૃતિક સંખ્યાઓનો ગણ}}\\
    \shoveleft{\Int = \{\ldots, -2, -1, 0, 1, 2, \ldots\}} \\
    \shoveright{\text{પૂર્ણાંકોનો ગણ}}\\
    \shoveleft{\Rat = \Setabs{\nicefrac{m}{n}}{m, n \in \Int\text{ અને }n \neq 0}}\\
    \shoveright{\text{સંમેય સંખ્યાઓનો ગણ}}\\
    \shoveleft{\Real = (-\infty, \infty)}\\
    \text{વાસ્તવિક સંખ્યાઓનો ગણ (સાતત્યક)}
\end{multline*}
આ બધા \emph{અનંત} ગણો છે, એટલે કે દરેકમાં અનંત સંખ્યામાં !!{element}s છે.

આ ગણોમાં આગળ વધતા જઈએ તેમ આપણી પાસે \emph{વધુ} સંખ્યાઓ ઉમેરાતી જાય છે.
હકીકતમાં, $\Nat \subseteq \Int \subseteq \Rat \subseteq \Real$ હોવું
સ્પષ્ટ થવું જોઈએ: દરેક પ્રાકૃતિક સંખ્યા પૂર્ણાંક છે; દરેક પૂર્ણાંક સંમેય
સંખ્યા છે; અને દરેક સંમેય સંખ્યા વાસ્તવિક સંખ્યા છે.
તેવી જ રીતે, $\Nat \subsetneq \Int \subsetneq \Rat$ હોવું પણ સ્પષ્ટ
થવું જોઈએ, કારણ કે $-1$ પૂર્ણાંક છે પરંતુ પ્રાકૃતિક સંખ્યા નથી,
અને $\nicefrac{1}{2}$ સંમેય છે પરંતુ પૂર્ણાંક નથી.
$\Rat \subsetneq \Real$, એટલે કે કેટલીક વાસ્તવિક સંખ્યાઓ સંમેય નથી,
તે એટલું સ્પષ્ટ નથી\oliflabeldef{sfr:arith:real:realline}{, પરંતુ
\olref[arith][real]{realline}માં આપણે આ મુદ્દા પર પાછા આવીશું}{}.

ક્યારેક આપણે ધન પૂર્ણાંકોનો ગણ $\PosInt = \{1, 2, 3, \dots\}$ અને
માત્ર પહેલી બે પ્રાકૃતિક સંખ્યાઓ ધરાવતો ગણ $\Bin = \{0, 1\}$ પણ વાપરીશું.
\end{ex}

\begin{tagblock}{compsci}
\begin{ex}[પ્રતીકશ્રેણીઓ]
બીજું રસપ્રદ ઉદાહરણ મૂળાક્ષરસમૂહ $A$ પરથી બનતી \emph{સાંત
પ્રતીકશ્રેણીઓ}નો ગણ $A^{*}$ છે: $A$ના ઘટકોનો કોઈ પણ સાંત અનુક્રમ
$A$ પરની પ્રતીકશ્રેણી છે. દરેક મૂળાક્ષરસમૂહ~$A$ માટે, $A$ પરની
પ્રતીકશ્રેણીઓમાં આપણે \emph{ખાલી પ્રતીકશ્રેણી $\Lambda$}નો સમાવેશ કરીએ છીએ.
દાખલા તરીકે,
\begin{multline*}
\Bin^*
=\{\Lambda,0,1,00,01,10,11,\\
000,001,010,011,100,101,110,111,0000,\ldots\}.
\end{multline*}
જો $x=x_{1}\ldots x_{n}\in A^{*}$ એ $A$ના $n$ “અક્ષરો”થી બનેલી
પ્રતીકશ્રેણી હોય, તો આપણે કહીએ કે તેની \emph{લંબાઈ}~$n$ છે
અને $\len{x}=n$ લખીએ.
\end{ex}
\end{tagblock}

\begin{ex}[અનંત અનુક્રમો]
કોઈ પણ ગણ $A$ માટે આપણે $A$ના !!{element}sના અનંત અનુક્રમોનો ગણ~$A^\omega$
પણ વિચારી શકીએ. અનંત અનુક્રમ $a_1a_2a_3a_4\dots$ એ વસ્તુઓની એક
દિશામાં અનંત સુધી ચાલતી યાદી છે, જેમાંની દરેક વસ્તુ $A$નો !!a{element} છે.
\end{ex}

\end{document}

